\renewcommand{\reportmonth}{Juni}
\renewcommand{\signaturedate}{2 Juli 2026}
\newpage
\TabPositions{4.5cm}

\begin{center}
    \textbf{LOG BOOK \\ MAGANG INDUSTRI}
\end{center}

\raggedright%
Program Studi           \tab: {\studyprogram} \\
Semester/Th. Akademik   \tab: {\semester} / {\academicperiod} \\
NIM/Nama Mahasiswa      \tab: {\authornim} / {\authorname} \\
Perusahaan              \tab: {\company} \\
Job Title               \tab: {\jobposition} \\
Departemen              \tab: {\jobdepartment} \\
Mentor                  \tab: {\mentor} \\
Bulan                   \tab: {\reportmonth} {\reportyear} \\
\vspace{12pt}

\begin{logtable}

    % =============================================================
    % Minggu Ke-21 (1 Juni -- 7 Juni 2026)
    % =============================================================

    \logentry{1.}
    {
        Selasa \linebreak 2 Jun 2026 \linebreak
        08:13--17:21 \linebreak % chktex 8
        WFO
    }{
        \begin{loglist}
            \item Melengkapi Task \#14696 (\textit{USER STORY}): Export file Dashboard Inventaris
            dengan format XLSX
            \item Menyelesaikan Task \#14740 (\textit{BUG}): Perbaikan \textit{sorting} kolom
            \texttt{jumlah\_data} saat \textit{group by} aktif
            \item Menyelesaikan Task \#14697 (\textit{USER STORY}): \textit{Group by} lebih dari
            1 kolom
        \end{loglist}
    }{
        \begin{logwork}
            \item (\#14696) Menambahkan fungsi \textit{export} ke format XLSX sebagai pelengkap dari fitur CSV; menambahkan pengecekan \texttt{hasActiveFilters} agar hasil ekspor sesuai filter aktif; menonaktifkan \textit{alternating row color} saat proses \textit{export}
            \item (\#14740) Memperbaiki logika pengurutan pada kolom \texttt{jumlah\_data} yang bermasalah ketika data sedang dikelompokkan
            \item (\#14697) Merapikan dan memfinalisasi implementasi \textit{multiple group by keys} pada \textit{DataView} dan Dashboard Inventaris, termasuk merapikan komponen \textit{chip} kelompok data
        \end{logwork}
    }

    \logentry{2.}
    {
        Rabu \linebreak 3 Jun 2026 \linebreak
        08:11--17:14 \linebreak % chktex 8
        WFO
    }{
        \begin{loglist}
            \item Mendapat dan menyelesaikan Task \#14858 (\textit{BUG}): Dropdown tidak lengkap
            menampilkan data di layar Riwayat Pengobatan
            \item Mendapat Task \#14965 (\textit{IMPROVEMENT}): Pembaruan
            \texttt{generateColumnDef} di \textit{DataView} agar otomatis mendeteksi format
            data
            \item Mendapat Task \#14857 (\textit{NEW}): Inisiasi pengerjaan \textit{alignment
                endpoint} Stok Opname \& Riwayat Belanja ke \textit{backend}
            \item Melanjutkan penanganan Task \#14374 (\textit{BUG}): Diagnosis Sekunder hilang
            di widget \textit{search chip} Flutter
        \end{loglist}
    }{
        \begin{logwork}
            \item (\#14858) Menemukan penyebab dropdown terkunci karena properti \texttt{prototypeItem} pada komponen \texttt{search\_chip}; menghapus properti tersebut agar tinggi \textit{chip} responsif secara dinamis
            \item (\#14965) Merancang fitur \textit{auto-detect} tipe data pada \texttt{generateColumnDef}: memformat otomatis tipe \textit{currency} ke 'Rp' dan tipe \textit{date} ke format tanggal Indonesia
            \item (\#14857) Mempelajari kebutuhan struktur \textit{endpoint} pada komponen Toraja untuk halaman terkait
        \end{logwork}
    }

    \logentry{3.}
    {
        Kamis \linebreak 4 Jun 2026 \linebreak
        12:08--17:11 \linebreak % chktex 8
        WFO *(terlambat 248 menit)*
    }{
        \begin{loglist}
            \item Menyelesaikan Task \#14374, \#14858, \#14965, \#14696, dan \#14697
        \end{loglist}
    }{
        \begin{logwork}
            \item (\#14374) Memperbaiki penanganan \textit{items provider} pada widget \texttt{search\_chip} Flutter agar data Diagnosis Sekunder tidak hilang saat AI Scribe melakukan pembaruan \textit{state}
            \item (\#14858) Melakukan \textit{commit} final perbaikan \texttt{search-chip} dan mengajukan \textit{merge request}
            \item (\#14965) Memfinalisasi implementasi otomatisasi \textit{data formatter} untuk \textit{CardItem} di komponen \textit{DataView}
            \item (\#14696 \& \#14697) Melakukan \textit{merge} cabang akhir untuk fitur \textit{export} XLSX dan \textit{multiple group by}, memastikan fungsionalitas berjalan lancar di \textit{environment staging}
        \end{logwork}
    }

    % =============================================================
    % Minggu Ke-22 (8 Juni -- 14 Juni 2026)
    % =============================================================

    \logentry{4.}
    {
        Senin \linebreak 8 Jun 2026 \linebreak
        08:03--17:14 \linebreak % chktex 8
        WFO
    }{
        \begin{loglist}
            \item Mendapat Task \#15085 (\textit{NEW}): Menghubungkan \textit{endpoint backend}
            ke halaman Details Belanja
            \item Membuat halaman \textit{landing page} baru aplikasi \textit{Stock Management}
            dan melakukan \textit{cleanup} komponen lama
        \end{loglist}
    }{
        \begin{logwork}
            \item (\#15085) Menganalisis struktur halaman Details Belanja, memetakan \textit{endpoint} Toraja, dan berkoordinasi dengan tim \textit{Backend} terkait format data \textit{response}
            \item Membangun halaman \textit{landing page} \textit{Stock Management} dengan
            \textit{routing} modul yang fleksibel (Riwayat Stok, Stok Opname, Dashboard
            Inventaris)
        \end{logwork}
    }

    \logentry{5.}
    {
        Selasa \linebreak 9 Jun 2026 \linebreak
        08:14--17:20 \linebreak % chktex 8
        WFO
    }{
        \begin{loglist}
            \item Melanjutkan Task \#14857 dan \#15085 (\textit{NEW}): membangun halaman Riwayat
            Belanja (\textit{List Belanja}) dan Details Belanja
            \item Membangun komponen \textit{PdfViewer} dan \textit{PdfModal} sebagai komponen
            pendukung untuk menampilkan dokumen faktur
        \end{loglist}
    }{
        \begin{logwork}
            \item (\#14857) Membangun struktur awal halaman Riwayat Belanja: \textit{coldef} tabel, \textit{response} \textit{schema}, \textit{hook}, \textit{service}, dan \textit{mapper} data
            \item (\#15085) Membangun struktur awal halaman Details Belanja: \textit{skeleton} \textit{loading}, \textit{response} \textit{schema}, \textit{hook} \texttt{use-get-details-belanja-data}, dan \textit{service} untuk \textit{fetch} data
            \item Membangun komponen \textit{PdfViewer} dan \textit{PdfModal} menggunakan
            \textit{library} \texttt{react-pdf-viewer} dan \texttt{pdfjs-dist} untuk
            menampilkan dokumen PDF (faktur pembelian) langsung di dalam \textit{modal}
        \end{logwork}
    }

    \logentry{6.}
    {
        Rabu \linebreak 10 Jun 2026 \linebreak
        08:12--17:19 \linebreak % chktex 8
        WFO
    }{
        \begin{loglist}
            \item Melanjutkan pengembangan halaman Details Belanja dan Riwayat Belanja
            \item Berkoordinasi dengan Tim \textit{Backend} terkait struktur \textit{endpoint}
        \end{loglist}
    }{
        \begin{logwork}
            \item (\#14857 \& \#15085) Menyusun daftar kebutuhan \textit{endpoint} untuk Details Belanja dan Riwayat Belanja (data faktur, status pembayaran, riwayat pengiriman) bersama Tim \textit{Backend}
            \item Menguji tampilan awal \textit{PdfViewer} dengan beberapa contoh dokumen PDF
            yang berbeda ukuran dan orientasi
        \end{logwork}
    }

    \logentry{7.}
    {
        Kamis \linebreak 11 Jun 2026 \linebreak
        07:58--17:18 \linebreak % chktex 8
        WFO
    }{
        \begin{loglist}
            \item Melanjutkan Task \#15085: \textit{alignment} Details Belanja ke
            \textit{backend}
            \item \textit{Alignment} Receiving List dan Details Rincian Receiving List ke
            \textit{backend}
        \end{loglist}
    }{
        \begin{logwork}
            \item (\#15085) Mere\textit{factor} \textit{DetailsBelanja}: memecah komponen \texttt{AlasanPembatalanModal} dan \texttt{PembayaranModal} menjadi \textit{subcomponent} tersendiri, menyesuaikan \textit{schema} dan \textit{utils}
            \item (\#15085) Melakukan \textit{alignment} Details Belanja ke \textit{endpoint backend}, menambahkan \textit{hook} \texttt{use-get-details-belanja-data}
            \item \textit{Alignment} Receiving List ke \textit{backend}: \textit{coldef}, \textit{response} \textit{schema}, \textit{hook} \texttt{use-get-receiving-list-table-data}, \textit{service}, dan \textit{mapper}; membersihkan kode yang tidak lagi digunakan
            \item Membangun dan meng-\textit{align} halaman Details Rincian Receiving List:
            \textit{skeleton}, \texttt{FakturReceiving}, \texttt{InformasiPurchaseOrder},
            \texttt{PenerimaanAccordion}, dan \texttt{RingkasanTransaksi}
        \end{logwork}
    }

    \logentry{8.}
    {
        Jumat \linebreak 12 Jun 2026 \linebreak
        08:14--17:16 \linebreak % chktex 8
        WFO
    }{
        \begin{loglist}
            \item Menyelesaikan perbaikan urutan data Penerimaan pada Details Rincian Receiving
            List
            \item Melanjutkan \textit{alignment} Riwayat Stok ke \textit{backend} mengikuti
            struktur baru
            \item Mendapat Task \#15229 (\textit{IMPROVEMENT}): pembuatan React \textit{Error
                Parser} dan \textit{Error Modal}
        \end{loglist}
    }{
        \begin{logwork}
            \item Membalik urutan tampilan \textit{accordion} Penerimaan pada Details Rincian
            Receiving List agar data terbaru muncul di bagian atas
            \item \textit{Alignment} ulang halaman Riwayat Stok: \textit{enum}, \textit{coldef}, \textit{response} \textit{schema}, \textit{hook}, \textit{service} riwayat, dan \textit{mapper} data mengikuti pola \textit{alignment} yang baru
            \item (\#15229) Menganalisis variasi struktur \textit{error response} dari \textit{backend} (\texttt{code}, \texttt{message}, \texttt{detail}) untuk merancang \textit{parser} dan \textit{modal} \textit{error} yang seragam di seluruh halaman
        \end{logwork}
    }

    % =============================================================
    % Minggu Ke-23 (15 Juni -- 21 Juni 2026)
    % =============================================================

    \logentry{9.}
    {
        Senin \linebreak 15 Jun 2026 \linebreak
        08:07--17:15 \linebreak % chktex 8
        WFO
    }{
        \begin{loglist}
            \item Merancang struktur komponen \textit{ErrorModal} dan fungsi \texttt{parseError}
        \end{loglist}
    }{
        \begin{logwork}
            \item (\#15229) Merancang struktur komponen \textit{ErrorModal} yang dapat menangani berbagai bentuk \texttt{error.detail} (\textit{string}, \textit{array}, maupun \textit{object}/\texttt{Record})
        \end{logwork}
    }

    \logentry{10.}
    {
        Rabu \linebreak 17 Jun 2026 \linebreak
        08:07--17:19 \linebreak % chktex 8
        WFO
    }{
        \begin{loglist}
            \item Melanjutkan pengembangan dan pengujian komponen \textit{ErrorModal}
            \item Mempersiapkan Task InfiniteSelect dan \textit{editable table} \textit{DataView}
        \end{loglist}
    }{
        \begin{logwork}
            \item (\#15229) Mengimplementasikan tampilan awal \textit{ErrorModal} dan menguji dengan beberapa skenario \textit{error} dari \textit{endpoint} yang sudah di-\textit{align}
            \item (\#15229) Menguji \textit{ErrorModal} lebih lanjut pada berbagai halaman yang sudah terhubung ke \textit{backend}
            \item Mempersiapkan rancangan komponen \textit{InfiniteSelect} (\textit{dropdown}
            dengan \textit{lazy loading}) untuk kebutuhan pemilihan distributor pada Tambah
            Pesanan dan Details Belanja
        \end{logwork}
    }

    \logentry{11.}
    {
        Kamis \linebreak 18 Jun 2026 \linebreak
        08:20--17:24 \linebreak % chktex 8
        WFO
    }{
        \begin{loglist}
            \item Mendapat Task \#15407 (\textit{NEW}): \textit{Align BE Endpoint} ke halaman
            Tambah Pesanan (Belanja)
            \item Mendapat Task \#15317 (\textit{NEW}): React
            \textit{InfiniteSelect}/\textit{Lazy Loader Dropdown}
            \item Mendapat Task \#14965 (\textit{IMPROVEMENT}) lanjutan: \textit{DataView
                editable table}
        \end{loglist}
    }{
        \begin{logwork}
            \item (\#15407) Membangun struktur awal \textit{alignment} halaman Tambah Pesanan: \texttt{page.tsx}, \textit{response} \textit{schema}, \textit{hook} \texttt{use-tambah-pesanan-page-data}, \textit{service}, dan \textit{mapper}
            \item (\#15317) Membangun komponen \textit{InfiniteSelect}: \textit{Select} dengan \textit{lazy loader} menggunakan \textit{React Query} untuk menampilkan data \textit{dropdown} secara bertahap (\textit{infinite scroll}), diterapkan pertama kali pada halaman Pilih Nomor Lot
        \end{logwork}
    }

    \logentry{12.}
    {
        Jumat \linebreak 19 Jun 2026 \linebreak
        08:07--17:31 \linebreak % chktex 8
        WFO
    }{
        \begin{loglist}
            \item Menyelesaikan Task \#15229 (React \textit{Error Parser} dan \textit{Error
                Modal})
            \item Melanjutkan Task \#14965: \textit{DataView editable table}
            \item Melanjutkan Task \#15317: \textit{InfiniteSelect} dan \textit{Create
                Distributor Modal}
        \end{loglist}
    }{
        \begin{logwork}
            \item (\#15229) Memfinalisasi \textit{ErrorModal}: menambahkan penanganan untuk \texttt{error.detail} berupa \texttt{code}, \texttt{message}, dan \texttt{detail} dengan format yang lebih baik
            \item (\#14965) Menambahkan fitur \textit{editable rows} pada \textit{DataView}: baris tabel dapat langsung diedit tanpa membuka \textit{modal} terpisah
            \item (\#15317) Menyempurnakan \textit{InfiniteSelect} agar dapat menyimpan \textit{state} berupa \textit{object} (bukan hanya \textit{string}), diperlukan untuk kebutuhan \textit{Create Distributor Modal}
            \item (\#15317) Membangun dan meng-\textit{align} \textit{Create Distributor Modal} menggunakan \textit{InfiniteSelect}; menghapus opsi input nomor telepon rumah (\textit{landline}) yang sudah tidak relevan; membatasi jumlah karakter input dan merapikan \textit{error handling}
        \end{logwork}
    }

    % =============================================================
    % Minggu Ke-24 (22 Juni -- 28 Juni 2026)
    % =============================================================

    \logentry{13.}
    {
        Senin \linebreak 22 Jun 2026 \linebreak
        08:05--17:15 \linebreak % chktex 8
        WFO
    }{
        \begin{loglist}
            \item Melanjutkan Task \#15407: \textit{alignment} Tambah Pesanan
        \end{loglist}
    }{
        \begin{logwork}
            \item (\#15407) Menentukan lebar kolom tabel, menambah dan menghapus \textit{field} wajib (\textit{required fields}), serta menambahkan \texttt{small\_unit\_name} pada \textit{label} pilihan varian produk
            \item (\#15407) Membuat tabel Tambah Pesanan menjadi \textit{editable}, menghubungkan ke \textit{endpoint} POST, dan menambahkan fitur tambah distributor baru langsung dari halaman ini
        \end{logwork}
    }

    \logentry{14.}
    {
        Selasa \linebreak 23 Jun 2026 \linebreak
        08:07--17:27 \linebreak % chktex 8
        WFO
    }{
        \begin{loglist}
            \item Menyelesaikan Task \#14857 dan \#15085: Riwayat Stok, List Belanja, Details
            Belanja, Receiving List, dan Pilih Nomor Lot
            \item Menyelesaikan Task \#15407 (\textit{Align BE Endpoint} Tambah Pesanan)
            \item Menyelesaikan Task \#15317 (\textit{InfiniteSelect})
            \item Migrasi halaman Item Faskes (Daftar Produk Item Cabang)
        \end{loglist}
    }{
        \begin{logwork}
            \item (\#15407) Menyelesaikan \textit{alignment} Tambah Pesanan agar dapat membuat penerimaan stok tanpa transaksi (\textit{create receiving stock without transaction}); menambahkan \textit{loading state} pada tombol \textit{submit}
            \item (\#14857) Menyempurnakan pemetaan \textit{reference}/\textit{activity type} pada Riwayat Stok serta merapikan perbaikan minor lainnya
            \item (\#14857) Menyelesaikan \textit{alignment} List Belanja: menambahkan penanganan pembayaran dan pembatalan pesanan
            \item (\#15085) Menyelesaikan \textit{alignment} Details Belanja: menambahkan penanganan pembayaran, pembatalan, serta \textit{upload} dan \textit{fetch} faktur
            \item Menyelesaikan \textit{alignment} Pilih Nomor Lot ke \textit{endpoint backend}
            yang baru dan Receiving List
            \item Memigrasi halaman Item Faskes (Daftar Produk Item Cabang) beserta seluruh
            subhalamannya (Layanan, Paket, Produk, Tindakan) ke struktur baru
            \item Menambahkan dukungan tampilan gambar pada \textit{PdfModal}
            \item Merapikan penanganan \textit{error} di berbagai halaman serta memperbaiki
            \texttt{parseError} agar kompatibel dengan struktur \textit{error} dari
            \textit{stock-v2}
            \item Membersihkan variabel lingkungan (\texttt{.env.local}) yang tidak lagi
            digunakan; memperbarui \textit{package} \textit{tanstack query} untuk kebutuhan
            \textit{lazy loader} \textit{InfiniteSelect}; mere\textit{factor}
            \textit{provider} aplikasi
        \end{logwork}
    }

    \logentry{15.}
    {
        Rabu \linebreak 24 Jun 2026 \linebreak
        08:05--17:25 \linebreak % chktex 8
        WFO
    }{
        \begin{loglist}
            \item Menambahkan fitur \textit{auto select} pada \textit{InfiniteSelect}
            \item Mendapat dan menyelesaikan Task \#15453 (\textit{IMPROVEMENT}): Flutter
            \textit{Disable Page Wrapper Padding}
        \end{loglist}
    }{
        \begin{logwork}
            \item Menambahkan fitur agar \textit{InfiniteSelect} otomatis memilih data pertama
            saat pertama kali dimuat, mengurangi langkah manual pengguna saat hanya
            tersedia satu pilihan
            \item (\#15453) Menambahkan opsi untuk menonaktifkan \textit{padding} bawaan pada \textit{Page Wrapper} di Flutter, digunakan pada layar tertentu yang membutuhkan tampilan penuh tanpa jarak tambahan
        \end{logwork}
    }

    \logentry{16.}
    {
        Kamis \linebreak 25 Jun 2026 \linebreak
        08:08--17:15 \linebreak % chktex 8
        WFO
    }{
        \begin{loglist}
            \item Mengimpor modul \textit{Packing List} (Pengemasan dan Pengeluaran) ke
            \textit{repository} baru
            \item Membangun komponen \textit{DevToolkit} untuk membantu proses
            \textit{development}
        \end{loglist}
    }{
        \begin{logwork}
            \item Mengimpor halaman \textit{Packing List} beserta seluruh subkomponennya
            (\textit{Cetak Resi, Detail Pesanan, Proses Kemas}, \textit{tab Dalam
                Pengiriman} dan \textit{Dibatalkan} termasuk versi \textit{online}) dari
            \textit{repository} lama
            \item Membangun komponen \textit{DevToolkit}: alat bantu pengembang untuk mempermudah
            proses \textit{debugging} dan pengujian selama \textit{development} berlangsung
        \end{logwork}
    }

    \logentry{17.}
    {
        Jumat \linebreak 26 Jun 2026 \linebreak
        08:09--18:49 \linebreak % chktex 8
        WFO
    }{
        \begin{loglist}
            \item Melakukan berbagai penyesuaian UI dan \textit{styling} pada halaman-halaman
            \textit{Stock Management}
            \item Menambahkan \textit{top horizontal scrollbar} pada \textit{DataView} serta
            menyesuaikan ukuran \textit{font}
            \item Melanjutkan \textit{alignment} halaman Add Receiving List
        \end{loglist}
    }{
        \begin{logwork}
            \item Melakukan penyesuaian UI dan \textit{styling} pada Details Belanja, Details
            Rincian Receiving List, List Belanja, Pilih Unit Gudang, Receiving List, dan
            Tambah Pesanan
            \item (\textit{DataView}) Menambahkan \textit{scrollbar} horizontal di bagian atas tabel agar navigasi lebih mudah pada tabel lebar; menyesuaikan seluruh ukuran \textit{font} dan elemen terkait; memperbarui \textit{styling scrollbar}
            \item Merapikan \textit{styling} \textit{FooterButton} dan \textit{InfiniteSelect}
            \item Memperbaiki tipe data \texttt{backToraja} dan \texttt{refreshToraja} pada
            \textit{redirect utils} agar sesuai dengan Flutter \textit{message listener}
            \item Melakukan \textit{alignment} halaman Add Receiving List ke \textit{backend}
            \item (\#15229) Menambahkan penanganan tambahan pada \textit{ErrorModal} untuk struktur \texttt{error.detail} berupa \texttt{Record}
        \end{logwork}
    }

    % =============================================================
    % Minggu Ke-25 (29 Juni -- 30 Juni 2026)
    % =============================================================

    \logentry{18.}
    {
        Senin \linebreak 29 Jun 2026 \linebreak
        08:21--17:30 \linebreak % chktex 8
        WFO
    }{
        \begin{loglist}
            \item Mendapat Task \#15410 (\textit{NEW}): \textit{Align BE Endpoint} ke Perpindahan
            Stok/\textit{Transfer Stock}
        \end{loglist}
    }{
        \begin{logwork}
            \item (\#15410) Membangun dan memulai \textit{alignment} halaman \textit{Transfer Stock}: \texttt{page.tsx}, \textit{TransferStock.tsx}, \textit{params}, \textit{response} \textit{schema}, \textit{coldef}, \textit{hook} \texttt{use-transfer-stock}, \textit{service}, serta \textit{subcomponent} \texttt{TransferStockReceipts} dan \texttt{TransferStockSubmissions}
        \end{logwork}
    }

    % =============================================================
    % Evaluasi Bulanan
    % =============================================================
    \logevaluate{Analisis dan Pembahasan}{
        Di awal bulan, pekerjaan berfokus pada penyelesaian sisa \textit{task} bulan Mei
        (\#14696, \#14740, \#14697) serta perbaikan \textit{bug} \textit{dropdown} pada
        Riwayat Pengobatan (\#14858) dan penyempurnaan \textit{auto-detect} format data
        pada \texttt{generateColumnDef} (\#14965).

        Sejak minggu kedua, fokus utama beralih ke \textit{alignment} modul Riwayat
        Belanja secara menyeluruh (Task \#14857 dan \#15085): halaman Riwayat Belanja,
        Details Belanja, Receiving List, Details Rincian Receiving List, dan Riwayat
        Stok dihubungkan ke \textit{endpoint backend} yang baru. Untuk mendukung
        kebutuhan modul ini, dibangun beberapa komponen generik baru yaitu
        \textit{PdfViewer}/\textit{PdfModal} (menampilkan faktur), \textit{ErrorModal}
        beserta \texttt{parseError} (Task \#15229, menyeragamkan penanganan
        \textit{error} dari \textit{backend}), dan \textit{InfiniteSelect} (Task
        \#15317, \textit{dropdown} dengan \textit{lazy loading} untuk data dalam jumlah
        besar seperti daftar distributor).

        Pada minggu keempat, \textit{alignment} diperluas ke halaman Tambah Pesanan
        (Task \#15407) dan migrasi halaman Item Faskes, ditandai dengan hari kerja yang
        cukup panjang pada 26 Juni untuk menyelesaikan seluruh rangkaian \textit{task}
        sebelum \textit{deadline}. Di sisa bulan, dilakukan penyempurnaan
        UI/\textit{styling} lintas halaman, penyelesaian Task \#15453 (Flutter), impor
        modul \textit{Packing List}, pembuatan komponen \textit{DevToolkit}, hingga
        dimulainya \textit{task} berikutnya yaitu \textit{alignment} \textit{Transfer
            Stock} (\#15410) yang akan dilanjutkan hingga awal Juli. \textit{Branch}
        penyelesaian Task \#14857 juga digabungkan ke \textit{branch}
        \texttt{refactor-stock-management} pada akhir bulan. }

    \logevaluate{Capaian yang Telah Diselesaikan}{
        \begin{itemize}[noitemsep, topsep=2pt, leftmargin=12pt, label=\textbullet]
            \item Menyelesaikan Task \#14696, \#14697, \#14740: fitur \textit{export} XLSX dan
                  \textit{multiple group by} pada \textit{DataView} dan Dashboard Inventaris.
            \item Menyelesaikan Task \#14858: perbaikan \textit{dropdown} Riwayat Pengobatan
                  serta \textit{React-Toraja Back and Refresh Button function}.
            \item Menyelesaikan Task \#14965: \textit{auto-detect} format data pada
                  \texttt{generateColumnDef} dan fitur \textit{editable table} pada
                  \textit{DataView}.
            \item Menyelesaikan Task \#14857 dan \#15085: \textit{alignment} Riwayat Belanja,
                  Details Belanja, Receiving List, Details Rincian Receiving List, dan Riwayat
                  Stok ke \textit{backend}.
            \item Menyelesaikan Task \#15229: komponen \textit{ErrorModal} dan
                  \texttt{parseError}.
            \item Menyelesaikan Task \#15317: komponen \textit{InfiniteSelect}/\textit{Lazy
                      Loader Dropdown}, termasuk \textit{Create Distributor Modal}.
            \item Menyelesaikan Task \#15407: \textit{alignment} halaman Tambah Pesanan.
            \item Menyelesaikan Task \#15453: Flutter \textit{Disable Page Wrapper Padding}.
            \item Membangun komponen \textit{PdfViewer}/\textit{PdfModal}, komponen
                  \textit{DevToolkit}, dan mengimpor modul \textit{Packing List}.
            \item Memigrasi halaman Item Faskes (Daftar Produk Item Cabang) ke struktur baru.
            \item Memulai Task \#15410 (\textit{Transfer Stock}) yang akan dilanjutkan di Juli.
        \end{itemize}
    }

    \logevaluate{Capaian yang Belum Diselesaikan}{
        \begin{itemize}[noitemsep, topsep=2pt, leftmargin=12pt, label=\textbullet]
            \item Task \#15410 (\textit{Align BE Endpoint Transfer Stock}) baru dimulai di akhir
                  bulan dan akan dilanjutkan hingga awal Juli.
        \end{itemize}
    }

    \logevaluate{Kendala}{
        \textit{Alignment} beberapa halaman secara paralel (Riwayat Belanja, Details
        Belanja, Receiving List, Tambah Pesanan) menuntut koordinasi yang cukup intensif
        dengan Tim \textit{Backend} karena setiap halaman memiliki kebutuhan
        \textit{endpoint} yang berbeda-beda namun saling terkait. Perbedaan struktur
        \textit{error response} antar \textit{endpoint} juga sempat menyulitkan proses
        penanganan \textit{error} di sisi \textit{front-end} sebelum \textit{ErrorModal}
        dan \texttt{parseError} selesai dibuat.
    }

    \logevaluate{Solusi}{
        Dibangun beberapa komponen generik (\textit{ErrorModal}, \textit{InfiniteSelect},
        \textit{PdfViewer}) untuk menstandardisasi cara penanganan \textit{error} dan
        pengambilan data lazim di seluruh halaman, sehingga tidak perlu ditangani secara
        berulang di setiap modul. Proses \textit{alignment} juga dikerjakan per halaman
        menggunakan \textit{branch} terpisah, yang kemudian digabungkan ke \textit{branch}
        \texttt{refactor-stock-management} setelah dipastikan stabil, guna menjaga
        kestabilan kode secara keseluruhan selama proses migrasi berlangsung.
    }

\end{logtable}
\begin{flushright}
    {\city}, {\signaturedate}\\
    Mengetahui,\\
    Mentor {\jenisMBKM}\\[2cm]
    {\mentor}\\
\end{flushright}